\documentclass[conference]{IEEEtran}
\usepackage{times}
\usepackage[numbers]{natbib}
\usepackage{multicol}
\usepackage[bookmarks=true]{hyperref}
\usepackage{url}
\usepackage{graphicx}
<<<<<<< HEAD
\usepackage{amsmath}
=======
>>>>>>> 1922a3f (Add 0528 project code reports and figures)

\pdfinfo{
   /Author (Haotian Wu, Yuhua Luo, Xuyang Yuan)
   /Title  (Style-Aware Adaptation for Language-Conditioned Humanoid Motion Generation)
   /Subject (Language-guided humanoid motion generation)
   /Keywords (MoConVQ; HumanML3D; LoRA; humanoid motion generation)
}

\title{Style-Aware Adaptation for Language-Conditioned Humanoid Motion Generation}
\author{\IEEEauthorblockN{Haotian Wu, Yuhua Luo, and Xuyang Yuan}}

\begin{document}
\maketitle

\begin{abstract}
<<<<<<< HEAD
This midterm report summarizes concrete progress on style-aware adaptation of MoConVQ for language-conditioned humanoid motion generation. Our project follows the language-guided motion generation direction: reproduce a recent physics-based text-to-motion baseline, identify a limitation in fine-grained style control, and adapt the language-conditioned generator while keeping the motion representation and controller frozen. Since the proposal, we have reproduced pretrained MoConVQ inference, generated ten BVH baseline rollouts, curated a split-preserving HumanML3D style-caption subset, diagnosed a format mismatch between the available HumanML3D parquet data and MoConVQ token supervision, and started a parameter-efficient LoRA distillation pipeline. Preliminary results include baseline motion snapshots, dataset statistics, and an initial 200-caption LoRA training curve. The remaining work is to complete adapted generation, compare baseline and LoRA BVH outputs, and report only measured metrics that are compatible with the generated files.
=======
This midterm report summarizes progress toward style-aware adaptation of MoConVQ for language-conditioned humanoid motion generation. We reproduced the pretrained MoConVQ inference pipeline, generated ten baseline BVH motions, built split-preserving HumanML3D style-subset tooling, and identified a key practical constraint: the available HumanML3D parquet export contains processed motion features rather than MoConVQ-compatible token targets. The next phase therefore uses parameter-efficient LoRA distillation as a faithful backup to the proposal's frozen-backbone adaptation plan.
>>>>>>> 1922a3f (Add 0528 project code reports and figures)
\end{abstract}

\IEEEpeerreviewmaketitle

<<<<<<< HEAD
\section{Problem Statement and Motivation}
The target problem is text-conditioned humanoid motion generation with explicit style control. Given a natural-language instruction such as ``walk forward slowly'' or ``move like doing tai chi,'' the system should generate a physically plausible humanoid motion that preserves both the requested action and the requested style. This is relevant to intelligent robotics because natural language is a convenient high-level command interface for humanoid robots, simulated characters, and embodied agents.

Our baseline is MoConVQ, which combines scalable discrete motion representations with physics-based motion control \cite{yao2024moconvq}. The baseline can generate BVH humanoid motions from prompts, but style-specific prompts are more difficult than generic action prompts because style is often expressed through motion speed, smoothness, arm range, vertical displacement, and body rotation rather than only through action category.

The proposal originally planned supervised fine-tuning on HumanML3D style motions. During implementation, we found that the available local HumanML3D export stores processed 263-D motion features and captions, while MoConVQ generation is supervised by its own discrete motion-token space. Directly treating the parquet features as MoConVQ tokens would be invalid. We therefore keep the original frozen-backbone goal but update the technical route: use the pretrained MoConVQ generator as a teacher to create pseudo-token targets, and train only LoRA updates in the language-conditioned generator.

\section{Related Work}
MoConVQ is the primary robotics baseline for this project. It learns a vector-quantized motion representation and uses physics-based control to track generated motion latents \cite{yao2024moconvq}. HumanML3D provides large-scale text-motion pairs and is widely used for text-driven motion generation \cite{guo2022humanml3d}. MDM generates motion using diffusion models over motion sequences \cite{tevet2023mdm}, while T2M-GPT demonstrates that discrete motion tokens are effective for text-to-motion generation \cite{zhang2023t2mgpt}. Our project follows the discrete-token and pretrained-generator direction, but focuses on a smaller parameter-efficient adaptation problem rather than retraining a motion representation model.

Our contribution is not a new motion backbone. Instead, we implement a project-specific adaptation layer around the existing MoConVQ codebase: style-caption filtering, baseline reproduction, pseudo-token distillation, LoRA injection, and BVH-level evaluation utilities. This makes the scope feasible within the course timeline and avoids training the expensive quantization/controller stack from scratch.

\section{Technical Approach}
The full pipeline is shown in Fig.~\ref{fig:pipeline}. A prompt is encoded by the text encoder, passed to the MoConVQ language-conditioned transformer, converted into discrete motion codes, decoded into motion latents, and tracked by the physics controller to produce BVH output. Our adaptation is restricted to the language-conditioned generator; the VQ codebooks, decoder, and controller remain frozen.
=======
\section{Problem and Motivation}
The goal is to generate humanoid motion from natural language while preserving both action semantics and requested style. MoConVQ is a strong robotics baseline because it uses discrete motion representations and physics-based tracking control \cite{yao2024moconvq}. However, style-specific prompts such as slow deliberate walking, tai chi-like movement, graceful dance, and heavy tired walking may be underrepresented in general text-to-motion training data.

\section{Related Work}
HumanML3D provides paired motion and text descriptions for text-driven motion generation \cite{guo2022humanml3d}. MDM models raw motion with diffusion \cite{tevet2023mdm}, while T2M-GPT demonstrates the effectiveness of discrete motion tokens for text-to-motion generation \cite{zhang2023t2mgpt}. Our project follows the discrete-token direction but keeps MoConVQ's motion representation and controller frozen.

\section{Technical Approach}
The intended proposal pipeline freezes the VQ representation, codebook, decoder, and controller, while adapting only the language-conditioned code generator.
>>>>>>> 1922a3f (Add 0528 project code reports and figures)

\begin{figure}[h]
\centering
\fbox{\begin{minipage}{0.94\linewidth}
\centering
Text prompt $\rightarrow$ text encoder $\rightarrow$ language transformer $\rightarrow$ discrete motion codes $\rightarrow$ MoConVQ decoder/controller $\rightarrow$ BVH humanoid motion
\end{minipage}}
<<<<<<< HEAD
\caption{Project pipeline. The midterm implementation has reproduced the baseline path and started LoRA adaptation inside the language-conditioned transformer.}
\label{fig:pipeline}
\end{figure}

\subsection{Style Data Curation}
We filter HumanML3D captions using style keyword groups for controlled motion, graceful smooth motion, dance-like motion, energetic dynamics, and heavy or awkward motion. Filtering is applied independently inside train, validation, and test splits to avoid leakage. The resulting small training split uses 200 style captions, with 40 validation and 40 test captions prepared for lightweight experiments.

\subsection{LoRA Distillation}
Let $x$ be a style caption, $z_{1:T}$ be the teacher-generated pseudo-token sequence, and $\theta$ be the frozen pretrained MoConVQ generator parameters. We add trainable LoRA parameters $\Delta\theta_{\phi}$ to selected linear layers and optimize cross-entropy over the teacher tokens:
\begin{equation}
\mathcal{L}(\phi) =
- \sum_{t=1}^{T}\log p_{\theta+\Delta\theta_{\phi}}(z_t \mid z_{<t}, x).
\end{equation}
Only the low-rank LoRA matrices are trainable. This matches the proposal's compute-efficient frozen-backbone spirit while avoiding unsupported HumanML3D-to-MoConVQ token conversion.

\section{Preliminary Experiments and Results}
\subsection{Baseline Reproduction}
The pretrained MoConVQ assets were installed and verified locally. We generated ten prompts covering generic and style-related motions, including walking, sitting, running in a circle, slow deliberate walking, graceful dance, tai chi-like movement, energetic jumping, and heavy tired walking. Fig.~\ref{fig:baseline} shows representative baseline BVH snapshots rendered from the generated files.

\begin{figure}[h]
\centering
\includegraphics[width=\linewidth]{../../outputs/figures/midterm_baseline_demo.png}
\caption{Baseline MoConVQ BVH snapshots generated from text prompts. These outputs verify that the pretrained inference, decoding, and BVH export pipeline are working before adaptation.}
\label{fig:baseline}
\end{figure}
=======
\caption{Project pipeline. Adaptation is restricted to the language-conditioned generator; the motion backbone remains frozen.}
\end{figure}

For data curation, we filter style-related captions independently within each HumanML3D split using five keyword groups: controlled motion, graceful smooth motion, expressive dance, energetic dynamics, and heavy/awkward motion. This prevents leakage across train, validation, and test sets.

\section{Midterm Progress}
Baseline reproduction is complete. The official pretrained MoConVQ assets are present locally, tokenization and decoding have been verified, and ten prompts have produced BVH files under \texttt{outputs/baseline/project\_prompts/}.

\begin{table}[h]
\centering
\caption{Implementation progress at midterm.}
\footnotesize
\begin{tabular}{p{0.30\linewidth}p{0.17\linewidth}p{0.31\linewidth}}
\hline
Component & Status & Evidence \\
\hline
Environment & Complete & CUDA PyTorch and local extensions verified \\
Baseline inference & Complete & 10 BVH files generated \\
Style filtering & Complete & Split-preserving parquet filter \\
Small splits & Complete & 200/40/40 style subset files \\
Fine-tuning & Planned & LoRA distillation selected after format diagnosis \\
\hline
\end{tabular}
\end{table}
>>>>>>> 1922a3f (Add 0528 project code reports and figures)

\begin{table}[h]
\centering
\caption{Style subset statistics from the available HumanML3D parquet export.}
\footnotesize
\begin{tabular}{lrr}
\hline
Split & Matched captions & Matched motion IDs \\
\hline
Train & 3332 & 2682 \\
Validation & 214 & 168 \\
Test & 586 & 486 \\
\hline
\end{tabular}
<<<<<<< HEAD
\label{tab:subset}
\end{table}

\subsection{Implementation Progress}
Table~\ref{tab:progress} summarizes implementation status at midterm. Compared with the proposal stage, the project has moved from planning to executable code: baseline generation, style-subset construction, small-split creation, and an initial LoRA distillation run are all present.

\begin{table}[h]
\centering
\caption{Implementation progress at midterm.}
\footnotesize
\begin{tabular}{p{0.30\linewidth}p{0.18\linewidth}p{0.32\linewidth}}
\hline
Component & Status & Evidence \\
\hline
Environment & Complete & CUDA PyTorch, local extensions, pretrained files verified \\
Baseline inference & Complete & 10 BVH files generated from prompts \\
Style filtering & Complete & Split-preserving HumanML3D parquet filter \\
Small splits & Complete & 200/40/40 style-caption splits \\
LoRA distillation & In progress & 200-caption training curve and checkpoints \\
Evaluation & Planned & BVH proxy metrics and qualitative comparison scripts prepared \\
\hline
\end{tabular}
\label{tab:progress}
\end{table}

\subsection{Preliminary Training Curve}
We ran an initial 200-caption LoRA distillation experiment to test whether the adapter can learn the teacher pseudo-token distribution. Fig.~\ref{fig:loss} shows that the epoch-mean cross-entropy decreases during this preliminary run. This result does not yet prove improved motion quality, but it verifies that the training path, LoRA parameter injection, cached pseudo-token targets, and checkpoint writing all function end-to-end.

\begin{figure}[h]
\centering
\includegraphics[width=\linewidth]{../../outputs/figures/lora_loss_curve.png}
\caption{Preliminary LoRA distillation loss on 200 style captions. The decreasing trend indicates that the adapter is learning teacher pseudo-token targets, but final motion-quality evaluation is still ongoing.}
\label{fig:loss}
\end{figure}

\section{Risks and Remaining Plan}
The main risk remains data compatibility. The available HumanML3D parquet features are not direct MoConVQ token labels, so official FID or R-Precision evaluation cannot be claimed until a verified feature-space bridge is implemented. To keep the project honest, the final report will separate three kinds of evidence: LoRA training loss, BVH kinematic proxy metrics, and qualitative baseline-versus-adapter comparison.

The remaining plan is as follows. First, generate adapted BVH outputs for the same prompts as the baseline. Second, compute BVH proxy metrics such as root speed, mean joint velocity, acceleration, jerk, vertical range, arm motion range, rotation amount, and duration. Third, render side-by-side baseline and adapted videos for qualitative inspection. If official HumanML3D evaluation remains incompatible with the available data format, we will explicitly report that limitation rather than fabricating benchmark numbers.

\section{Conclusion}
At midterm, the project has moved beyond the proposal stage. We have reproduced MoConVQ baseline generation, curated style-focused HumanML3D metadata, diagnosed the central supervision-format risk, and started a LoRA distillation adaptation pipeline with measured training loss. The project is therefore technically active and on track for a final submission based on measured training, generated BVH outputs, and transparent limitations.
=======
\end{table}

\section{Risk and Next Step}
The main technical risk is format mismatch. The available HumanML3D parquet export contains processed 263-D features, while MoConVQ supervised fine-tuning requires discrete token sequences in its own representation. A direct conversion is nontrivial and could introduce invalid supervision.

The final-stage plan is therefore to use a parameter-efficient backup: pretrained MoConVQ serves as a teacher to generate pseudo token targets for style captions, and a student trains only LoRA updates in the language-conditioned transformer. We will report only measured loss, generated BVH outputs, and qualitative comparisons unless a compatible FID/R-Precision evaluator is implemented.
>>>>>>> 1922a3f (Add 0528 project code reports and figures)

\bibliographystyle{plainnat}
\bibliography{../shared/references}

\end{document}
